\centerline{\bf \Huge }
\centerline{\bf \Huge Гомотетия.}

\begin{flushright}
        \scriptsize{Глооооорп!}
\end{flushright}

\problem{0.1}\textbf{Базовые гомотетии}
\begin{enumerate}
    \item Любые две окружности переводятся друг в друга гомотетией (вне зависимости от их взаимного расположения.)
    \item Касающиеся окружности переводятся друга в друга гомотетией с центром в точке касания
    \item Любые подобные треугольники переводятся друг в друга посредством поворотной гомотетии
    \item Любой треугольник гомотетичен своему серединному треугольнику
    \item Любые две параллельные прямые (даже если они совпадают) гомотетичны
    \item Любые два параллельных отрезка гомотетичны.
    \item Гомотетия переводит точку касания в точку касания
\end{enumerate}

\problem{1}
С помощью задачи $0.1.(6)$ докажите \textit{замечательное свойство тpaпеции}: точка пересечения диагоналей, точка пересечения продолжений боковых сторон и середины оснований любой трапеции лежат на одной прямой.


\problem{2}
Даны две параллельные прямые $l_1$ и $l_2$. Окружности $\omega_1$ и $\omega_2$ касаются друг друга внешним образом в точке $B$, окружность $\omega_1$ касается $l_1$ в точке $A$ и $\omega_2$ касается $l_2$ в точке $C$. Докажите, что точки $A,B \text{ и } C$ лежат на одной прямой.

\problem{3}
Внутри квадрата $ABCD$ взята точка $M$. Докажите, что точки пересечения медиан треугольников $AMB, BMC, CMD, DMA$ образуют квадрат.

\problem{4}
На основаниях $BC$ и $AD$ трапеции $ABCD$ во внешнюю сторону построены равносторонние треугольники $BCK$ и $ADL$. Докажите, что $K,L$ и точка пересечения диагоналей лежат на одной прямой.

\problem{5}
Две окружности касаются в точке $К$. Прямая, проходящая через точку $К$, пересекает эти окружности в точках $A$ и $B$. Докажите, что касательные к окружностям, проведенные через точки $A$ и $B$, параллельны.

\newpage

\hspace{3mm}

\problem{6}
Две окружности касаются в точке $K$. Через точку $K$ проведены две прямые, пересекающие первую окружность в точках $A$ и $B$, вторую - в точках $C$ и $D$. Докажите, что $A B \| C D$.

\problem{7}
На сторонах $A B$ и $A C$ треугольника $A B C$ взяты соответственно точки $M$ и $N$, причём $M N \| B C$. На отрезке МN взята точка $P$, причём $M P=\frac{1}{3} M N$. Прямая $A P$ пересекает сторону $B C$ в точке $Q$. Докажите, что $B Q=\frac{1}{3} B C$.

\problem{8}
а) Вписанная окружность треугольника $A B C$ касается стороны $A C$ в точке $D, D M$ - её диаметр. Прямая $B M$ пересекает сторону $A C$ в точке $K$. Докажите, что $A K=D C$.

\noindent
б) В окружности проведены перпендикулярные диаметры $A B$ и $C D$. Из точки $M$, лежащей вне окружности, проведены касательные к окружности, пересекающие прямую $A B$ в точках $E$ и $H$, а также прямые $M C$ и $M D$, пересекающие прямую $A B$ в точках $F$ и $K$. Докажите, что $E F=K H$.